% =============================================================================
% CHAPTER 1 - INTRODUCTION
% =============================================================================

\chapter{INTRODUCTION}

\ifpdf
	\graphicspath{{introduction/figures/PNG}{introduction/figures/PDF}{introduction/figures}}
\else
	\graphicspath{{introduction/figures/EPS}{introduction/figures}}
\fi

\section{Background}
The background section of your work

\section{Problem Statement}
The problem statement of your work
% sample
% The problem of hallucination in large language models (LLMs) has emerged as a critical barrier to their reliable deployment in knowledge-intensive applications. Despite their impressive fluency and generalisation capabilities, LLMs often generate outputs that are factually incorrect or inconsistent with external knowledge sources, a phenomenon known as hallucination \citep{ji2023survey,huang2023survey}. This issue is particularly acute in tasks such as knowledge graph question answering (KGQA), where the model must generate answers that are not only linguistically coherent but also structurally faithful to an underlying knowledge graph \citep{lan2022survey}.

\section{Aim and Objectives}

The aim of this project is to ........

To achieve this aim, the following objectives have been set:......

\section{Tools and Facilities Used}

\subsection{Software Tools}

\subsection{Facilities}

\section{Methods Used}

The project employed the following methods:

\section{Scope of Work}



\section{Organization of Work}


